\chapter{Elementi di calcolo tensoriale}
Avendo appurato le strutture di varietà e fibrati tangenti e cotangenti, si può introdurre il formalismo che tratterà le grandezze.
% ORSO: specifica le grandezze
\section{Tensori}
    Si introducono i tensori come funzionali multilineari (lineari su tutte le entrate) su uno spazio vettoriale e sul suo duale.
    \begin{definition}
    \footnote{
    in questa definizione si farà uso della notazione "$\mathbb{V} \times \cdots \times \mathbb{V}$" sottolineata dalla graffa che esprime il numero (k) di spazi moltiplicati cartesianamente, solo per chiarire esplicitamente l'idea di prodotto cartesiano ripetuto tante volte quanto il rango del tensore; da qui in avanti verrà usata la notazione più stringata $\bigtimes^k \mathbb{V}$
    }
        Sia $\mathbb{V}$ spazio vettoriale reale finito-dimensionale, allora:
        \begin{itemize}
            \item un \textbf{tensore covariante} di \textbf{rango} $k$ è una mappa multilineare \\
            $\tau : \underbrace{\mathbb{V} \times \cdots \times \mathbb{V}}_{k \mathrm{-volte}} \longrightarrow \mathbb{R},$ l'insieme di tutti i tensori covarianti è indicato da $T^k\mathbb{V};$
            \item un \textbf{tensore controvariante} di \textbf{rango} $p$ è una mappa multilineare \\ 
            $Y : \underbrace{\mathbb{V^*} \times \cdots \times \mathbb{V^*}}_{p \mathrm{-volte}} \longrightarrow \mathbb{R},$ l'insieme di tutti i tensori controvarianti è \mbox{indicato da $T_p\mathbb{V};$}
            \item un \textbf{tensore misto} di \textbf{rango} $(k, p)$ è una mappa multilineare \\
            $ F : \underbrace{\mathbb{V} \times \cdots \times \mathbb{V}}_{k \mathrm{-volte}} \times \underbrace{\mathbb{V^*} \times \cdots \times \mathbb{V^*}}_{p \mathrm{-volte}} \longrightarrow \mathbb{R},$ l'insieme di tutti i tensori misti è \mbox{indicato da $T^k_p\mathbb{V}$.}
        \end{itemize}
    \end{definition}
    
    Notare che $T^k\mathbb{V}, T^k\mathbb{V}$ e $T^k_p\mathbb{V}$ possono essere dotati della struttura di spazio vettoriale reale.
    Facciamo un esempio per chiarire: siano $\alpha, \beta \in \mathbb{V}^*$, definiamo $\tau_{\alpha, \beta} : \mathbb{V} \times \mathbb{V} \longrightarrow \mathbb{R}$ tale che $(\mathbf{v}_1, \mathbf{v}_2) \mapsto \alpha(\mathbf{v}_1) \beta(\mathbf{v}_2)$, allora ovviamente $\tau_{\alpha, \beta} \in T^2\mathbb{V}$. Inoltre, dato $a \in \mathbb{R},$ avremo che sarà ovviamente multilineare:
    \begin{align}
     &a \tau_{\alpha, \beta} = \tau_{a \alpha, \beta} = \tau_{\alpha, a \beta}
    &\begin{cases}
        \tau_{\alpha + \alpha', \beta} = \tau_{\alpha, \beta} + \tau_{\alpha', \beta} \\
        \tau_{\alpha, \beta + \beta'} = \tau_{\alpha, \beta} + \tau_{\alpha, \beta'}
    \end{cases}
    \end{align}
    Siano, ora, $\tau \in T^k\mathbb{V}$ e $\rho \in T^m\mathbb{V}$; definiamo un elemento di $T^{k+m}\mathbb{V}$ come
    \[
    (\tau \otimes \rho)(\mathbf{v}_1, \dots, \mathbf{v}_{k+m}) = \tau(\mathbf{v}_1, \dots, \mathbf{v}_{k}) \rho(\mathbf{v}_{k+1}, \dots, \mathbf{v}_{m})
    \]
    da cui è immediato vedere che il \textbf{prodotto tensore} è \emph{associativo}: $(\tau \otimes \rho) \otimes \sigma = \tau \otimes (\rho \otimes \sigma)$.

\section{Tensori e pseudo-simmetrie}
    Prendiamo ora in considerazione una classe molto importante di tensori: sia $\sigma \in S_k$ una permutazione di $k$ oggetti e $\operatorname{sgn}(\sigma)$ il segno della permutazione;
    \begin{definition}
        $\rho \in T^k\mathbb{V}$ è detto \textbf{simmetrico} se
        \[
        \rho(\mathbf{v}_1, \dots, \mathbf{v}_k) = \rho(\sigma(\mathbf{v}_1), \dots, \sigma(\mathbf{v}_k)),
        \]oppure, se 
        \[
        \rho(\mathbf{v}_1, \dots, \mathbf{v}_k) = \operatorname{sgn}(\sigma)\rho(\sigma(\mathbf{v}_1), \dots, \sigma(\mathbf{v}_k))
        \]
        sarà detto \textbf{a simmetria alterna}.
        Gli insiemi di tensori simmetrici e a simmetria alterna saranno rispettivamente indicati da $\Sigma^k\mathbb{R}$ e $\Lambda^k\mathbb{V}$.
    \end{definition}
    Dato un generico tensore covariante $\rho \in T^k\mathbb{V}$, esistono due mappe per ottenerne uno simmetrico o alternante simmetrico:
    \begin{align}
    \operatorname{Symm} : T^k\mathbb{V} \longrightarrow \Sigma^k\mathbb{R}, \qquad \operatorname{Symm}(\rho)(\mathbf{v}_1, \dots, \mathbf{v}_k) = \frac{1}{k!}\sum_{\sigma \in S_k}\rho(\sigma(\mathbf{v}_1), \dots, \sigma(\mathbf{v}_k)) \\
    \operatorname{Alt} : T^k\mathbb{V} \longrightarrow \Lambda^k\mathbb{V}, \qquad \operatorname{Alt}(\rho)(\mathbf{v}_1, \dots, \mathbf{v}_k) = \frac{1}{k!}\sum_{\sigma \in S_k}\operatorname{sgn}(\sigma)\rho(\sigma(\mathbf{v}_1), \dots, \sigma(\mathbf{v}_k))
    \end{align}
    \begin{definition}
    Dati due tensori a simmetria alterna $\tau \in \Lambda^k\mathbb{V}$ e $\rho \in \Lambda^m\mathbb{V}$, possiamo definire il \textbf{prodotto wedge} come
    \begin{equation}
        \tau \wedge \rho = \frac{(k+m)!}{k! \; m!}\operatorname{Alt}(\tau \otimes \rho)
    \end{equation}
    \end{definition}
    Sostituendo la definizione di $\operatorname{Alt}(\tau \otimes \rho)$:
    \[
    \tau \wedge \rho = \frac{1}{k! \; m!} \sum_{\sigma \in S_k}\operatorname{sgn}(\sigma) \; \tau(\sigma(\mathbf{v}_1), \dots, \sigma(\mathbf{v}_k)) \; \rho(\sigma(\mathbf{v}_{k+1}), \dots, \sigma(\mathbf{v}_{k+m}))
    \]
    \begin{property}
        Dati $\rho, \tau, \pi$ rispettivamente $\in \Lambda^{k_\rho, k_\tau, k_\pi} \mathbb{V}$, valgono:
        \begin{itemize}
            \item associativa $\qquad \quad \qquad (\rho \wedge \tau) \wedge \pi = \rho \wedge (\tau \wedge \pi)$ 
            \item pseudo-simmetria $\qquad \: \tau \wedge \rho = (-1)^{km}(\rho \wedge \tau)$ 
        \end{itemize}
    \end{property}
    Si enunciano, ora, una serie di risultati che specificano le caratteristiche dei tensori e gli insiemi di cui essi fanno parte.
    \begin{proposition}
        Siano $\{ \mathbf{b}_i\}$ e $\{ \bm{\beta}_i\}$ basi duali rispettivamente per $\mathbf{V}$ e $\mathbf{V}^*$:
        \begin{itemize}
            \item allora il set di tutti i tensori covarianti di rango $k$ della forma $\bm{\beta}^{i_1} \otimes \cdots \otimes \bm{\beta}^{i_k}$ forma una base per $T^k\mathbb{V}$;
            \item l'insieme $\{ \bm{\beta}_{i_1} \wedge \cdots \wedge \bm{\beta}_{i_k} \}$, con $i_1 < i_2 < \dots < i_k$, è base per $\Lambda^k\mathbb{V}$
        \end{itemize}
    \end{proposition}
    Potrà essere utile scrivere un elemento di $\Lambda^k\mathbb{V}$ come
    \[
    \omega = \frac{1}{k!}\omega_{i_1, \dots, i_k} \bm{\beta}_{i_1} \wedge \cdots \wedge \bm{\beta}_{i_k}
    \]
    in cui non è più importante l'ordinamento degli indici richiesto alla proposizione precedente, a causa del quale compare il fattore $1/k!$ e i coefficienti $\omega_{i_1, \dots, i_k}$ sono completamente \emph{antisimmetrici}. La struttura che abbiamo dato fin'ora di tensori simmetrici e simmetrici alternanti è una generalizzazione delle matrici simmetriche e antisimmetriche.

    Consideriamo ora un esempio per chiarire: sia $\Lambda^2\mathbb{R}^3$ con $\{ \mathbf{e}_i \}$ base canonica e $\{ \bm{\epsilon}^j \}$ base duale. Sia $\omega \in \Lambda^2\mathbb{R}^3$
    \[
    \omega = \omega_{12} \bm{\epsilon}^1 \wedge \bm{\epsilon}^2 + \omega_{13} \bm{\epsilon}^1 \wedge \bm{\epsilon}^3 + \omega_{23} \bm{\epsilon}^2 \wedge \bm{\epsilon}^3
    \]
    ora usiamo l'antisimmetria dei coefficienti $\omega_{ij}$:
    \[
    \omega_{ij} = -\omega_{ji} \qquad \Longrightarrow \qquad \omega = \frac{1}{2}\omega_{ij} \bm{\epsilon}^i \wedge \bm{\epsilon}^j
    \]
    inoltre,
    \[
    \omega(\mathbf{e}_1, \mathbf{e}_2) = 2 \omega_{12} \operatorname{Alt}(\bm{\epsilon}^1 \otimes \bm{\epsilon}^2) (\mathbf{e}_1, \mathbf{e}_2) = \omega_{12}
    \]

    \section{Prodotto tensore}
    Dato l'esempio della sezione precedente, si potrebbe pensare di poter scrivere
    \[
    T^k\mathbb{V} = \bigotimes^k \mathbb{V}^*
    \]
    ma stiamo tralasciando una particolarità importante: riprendiamo l'esempio del tensore $\tau_{\alpha, \beta} : \mathbb{V} \times \mathbb{V} \longrightarrow \mathbb{R}$ tale che $(\mathbf{v}_1, \mathbf{v}_2) \mapsto \alpha(\mathbf{v}_1) \beta(\mathbf{v}_2)$. Questo è un esempio di tensore \emph{decomponibile}:
    \[
    \alpha \otimes \beta \coloneqq \alpha_i \beta_j \bm{\beta}^i \otimes \bm{\beta}^j
    \]
    cioè di un tensore scrivibile come combinazione lineare dei prodotti tensoriali tra covettori diversi della base duale.
    Ma non tutti i tensori sono decomponibili, e.g.
    \[
    \bm{\epsilon}^1 \otimes \bm{\epsilon}^2 + 2 \bm{\epsilon}^2 \otimes \bm{\epsilon}^2 \quad \in T^2\mathbb{R}^2
    \]
    non è rappresentabile nella forma $\alpha \otimes \beta$ con $\alpha, \beta \in (\mathbb{R}^2)^*$.

    Diamo ora una trattazione più formale: sia $S$ un generico insieme di elementi; cerchiamo di creare una sorta di combinazione lineare a prescindere dalla natura dei suoi elementi.
    Prendiamo le \emph{combinazioni lineari formali} di elementi di $S$, i.e. $\sum_i a_i x_i$ con $a_i \in \mathbb{R}, x_i \in S$.
    Formalmente diremo che una combinazione lineare formale è una funzione $f : S \longrightarrow \mathbb{R}$ nulla dappertutto a meno di un sottoinsieme \emph{finito} di $S$.
    Identifichiamo con $x^*$ la funzione nulla su tutto $S - \{ x_i \}$ con $x_i$ insieme dei punti di $S$ in cui $f \neq 0$, e poniamo $x^* (x_i)= 1$ vale allora che $f= \sum_i a_i x_i^*$.
    La struttura qua descritta rappresenta la nozione di \emph{spazio vettoriale libero}, detto $\mathbb{R} \langle S \rangle$, ossia l'insieme di tutte le combinazioni lineari formali da elementi di un insieme $S$.

    Adesso possiamo dare una nozione di prodotto tensore tra spazi più rigorosa.
    Siano $\mathbb{V}$ e $\mathbb{W}$ spazi vettoriali finito-dimensionali e ne costruiamo un altro fatto dalle combinazioni lineari formali di oggetti del tipo $(\alpha, \mu) \in \mathbb{V}^* \times \mathbb{W}^*$, visti come mappe multilineari $\mathbb{V}^* \times \mathbb{W}^* \longrightarrow \mathbb{R}$.
    Sia $I$ il sottospazio \emph{spannato} da tutti gli elementi del tipo
    \begin{align*}
        a(\alpha, \mu) - (&a \alpha, \mu) \\
        a(\alpha, \mu) - (&\alpha, a \mu) \\
        (\alpha + \alpha', \mu) - (\alpha, &\mu) - (\alpha', \mu) \\
        (\alpha, \mu + \mu') - (\alpha, &\mu) - (\alpha, \mu')
    \end{align*}
    Notare che questi elementi a priori non sono uguali fra loro, perché non è garantita la bilinearità non sapendo la natura degli oggetti presenti in negli spazi duali.
    \begin{definition}
        Ora definiamo il \textbf{prodotto tensore} tra $\mathbb{V}^*$ e $\mathbb{W}^*$ come
        \begin{equation}
            \mathbb{V}^* \otimes \mathbb{W}^* \coloneqq \mathbb{R}\left \langle \mathbb{V}^* \times \mathbb{W}^* \right \rangle / I
        \end{equation}
    \end{definition}
    Dove $A/B$ è l'operazione insiemistica del quoziente con la quale si chiede di creare un insieme di elementi di $A$ in cui gli elementi di $B$ sono considerati \emph{nulli}.
    Infatti, abbiamo chiesto che il prodotto tensore sia lo spazio vettoriale libero di elementi di $\mathbb{V}^* \times \mathbb{W}^*$, prendendo nulli gli elementi che spannano $I$: quindi, otteniamo automaticamente le formule di bilinearità
    \begin{align}
        a(\alpha \otimes \mu) = (a \alpha) \otimes \mu = \alpha \otimes (a \beta) \\
        \begin{cases}
            \alpha \otimes \mu + \alpha' \otimes \mu = (\alpha + \alpha') \otimes \mu \\
            \alpha \otimes \mu + \alpha \otimes \mu' = \alpha \otimes (\mu + \mu') \\
        \end{cases}
    \end{align}
    \begin{corollary}
        Sia $\mathbb{V}$ spazio vettoriale finito-dimensionale, allora
        \[
        T^k_r \mathbb{V} = \left(\bigotimes^k \mathbb{V}^*\right) \otimes \left(\bigotimes^r \mathbb{V}\right)
        \]
    \end{corollary}
    % ORSO
    % Alcuni esempi banali per capire questa struttura appena descritta sono:
    % \begin{itemize}
        % \item $T^0_0 \mathbb{V} = \mathbb{R}$;
        % \item $T^1 \mathbb{V} = \mathbb{V}^*$;
        % \item $T_1 \mathbb{V} = \mathbb{V}$;
        % \item
    % \end{itemize}
    Si enunciano, ora, un paio di risultati che mostrano come uno spazio tensoriale abbia proprietà di isomorfismo molto precise.
    In particolare, si può trovare un \emph{isomorfismo canonico}, detto anche \emph{naturale}, cioè un isomorfismo che relaziona esattamente gli oggetti di due spazi distinti, ma indipendentemente dalla scelta delle basi che li descrivono.
    Un isomorfismo generico cambierebbe la sua forma esplicita al variare delle basi dei due spazi, quello canonico no, per cui risulta che i due spazi sono intrinsecamente identici.
    Da un punto di vista simbolico, si passa da $\simeq$ a $=$ tra i due spazi.
    \begin{proposition}
        \begin{enumerate} Isomorfismi canonici tra spazi tensoriali e forme bi- o multi-lineari:
            \item Lo spazio vettoriale $\mathbb{V}^* \otimes \mathbb{W}^*$ è \textbf{canonicamente isomorfo} a $\operatorname{Bil}(\mathbb{V}, \mathbb{W})$, spazio delle funzioni bilineari $\mathbb{V} \times \mathbb{W} \longrightarrow \mathbb{R}$;
            \item $T^k_{p+1} \mathbb{V}$ è \textbf{canonicamente isomorfo} allo spazio delle mappe multilineari
            \[
            \left(\bigtimes^k \mathbb{V}\right) \times \left(\bigtimes^p \mathbb{V}^*\right) \longrightarrow \mathbb{V}
            \]
        \end{enumerate}
    \end{proposition}
    La corrispondenza nel rango tra $T^k_{p+1}$ e le mappe multilineari definite sopra può sembrare a prima vista errata.
    In realtà bisogna prestare attenzione all'output delle funzioni: esse restituiscono un vettore, mentre un tensore è definito come una mappa che mappa nel campo dei reali.
    In questo modo stiamo creando una corrispondenza tra quelle mappe e uno spazio tensoriale che contiene un covettore in più (il "+1") che manderà l'output delle mappe in uno scalare.
    Per esempio un endomorfismo su uno spazio vettoriale è una matrice $\mathbb{V} \longrightarrow \mathbb{V}$.
    Per renderla un tensore dovrà restituire uno scalare, quindi apparterrà allo spazio tensoriale $T^1_1\mathbb{V}$, i.e. avrà rango totale 2.
    
    \section{Fibrati tensoriali}
    Cerchiamo ora di rendere le strutture di spazi tensoriali covarianti, controvarianti e misti, dei fibrati \emph{tensoriali}.
    \begin{definition}
        Sia $\mathcal{M}$ una varietà differenziabile e $\Gamma(\cdot)$ la sezione locale del fibrato vettoriale in input:
        \begin{enumerate}
            \item il \textbf{fibrato tensoriale covariante} di rango $k$ sarà $T^k\mathcal{M} \coloneqq \bigsqcup_{p \in \mathcal{M}} T^k(T_p\mathcal{M})$, il corrispondente \textbf{campo tensoriale covariante} sarà $\mathcal{T}^k(\mathcal{M}) \coloneqq \Gamma (T^k\mathcal{M})$;
            \item il \textbf{fibrato tensoriale controvariante} di rango $r$ sarà $T_r\mathcal{M} \coloneqq \bigsqcup_{p \in \mathcal{M}} T^k(T_p\mathcal{M})$, il corrispondente \textbf{campo tensoriale controvariante} sarà $\mathcal{T}_r(\mathcal{M}) \coloneqq \Gamma (T_r\mathcal{M})$;
            \item il \textbf{fibrato tensoriale misto} di rango $(k,r)$ sarà $T_r^k\mathcal{M} \coloneqq \bigsqcup_{p \in \mathcal{M}} T_r^k(T_p\mathcal{M})$, il corrispondente \textbf{campo tensoriale misto} sarà $\mathcal{T}_r^k(\mathcal{M}) \coloneqq \Gamma (T_r^k\mathcal{M})$.
        \end{enumerate}
    \end{definition}
    Diamo adesso l'interpretazione \emph{pratica} assumendo una carta di coordinate, descritte dalle funzioni $x^i$, sulla varietà.
    Un qualsiasi tensore $T \in \mathcal{T}^k_r(\mathcal{M})$ sarà rappresentabile localmente come
    \begin{equation}
        T = 
        T^{j_1, \dots, j_r}_{i_1, \dots, i_k} \odif{x^{i_1}} \otimes \cdots \otimes \odif{x^{i_k}} \otimes \partial_{j_1} \otimes \cdots \otimes \partial_{j_r} =
        T^{j_1, \dots, j_r}_{i_1, \dots, i_k} \left( \bigotimes_{m=1}^k \odif{x^{i_m}} \right) \otimes \left( \bigotimes_{n=1}^r \partial_{j_n} \right)
    \end{equation}
    dove $T^{j_1, \dots, j_r}_{i_1, \dots, i_k}$ sono le funzioni lisce che rappresentano le coordinate del tensore $T$ rispetto alle basi duali degli spazi tangente e cotangente $\{ \odif{x^i} \}_{i_1}^{i_k}$ e $\{ \partial_j \}_{j_1}^{j_r}$.
    Diremo che $T_p \in T_r^k(T_p \mathcal{M})$ è il valore del campo tensoriale nel punto $p \in \mathcal{M}$, definito come
    \begin{equation}
        T_p(X_{1_p}, \dots, X_{2_p}, \omega_{1_p}, \dots, \omega_{2_p}) = T(X_{1}, \dots, X_{k}, \omega_{1}, \dots, \omega_{2}) (p)
    \end{equation}
    Vediamolo di nuovo in coordinate, più esplicitamente: sia ora $S \in \mathcal{T}_1^1(\mathcal{M})$; nella carta di coordinate sarà $S= S^j_i \odif{x^i} \otimes \partial_j$.
    Facendolo agire su un vettore $\mathbf{v}= v^i \partial_i$ e un covettore $\omega = \omega_i \odif{x^i}$, otterremo
    \begin{align*}
    S(\mathbf{v}, \omega) &= S^j_i v^k \omega_l \odif{x^i} (\partial_k) \odif{x^l} (\partial_j) \\
    \text{usando } \odif{x^\sigma}(\partial_\rho)=\delta_\rho^\sigma \quad \longrightarrow \quad &= S_i^j v^i \omega_j
    \end{align*}
    Risulta ovvio che la mappa trovata sia lineare e $\mathcal{C}^\infty (\mathcal{M})$.
    Questo risultato è generalizzabile a un generico campo tensoriale, per cui un campo tensoriale (1, 1) è una mappa multilineare $\mathcal{C}^\infty (\mathcal{M})$ da $\mathcal{X}(\mathcal{M}) \times \Omega^1(\mathcal{M}) \longrightarrow \mathcal{C}^\infty (\mathcal{M})$.
    Possiamo dare un risultato ancora più generale:
    \begin{lemma}
        Una mappa che sia o $\left( \bigtimes^k \mathcal{X}(\mathcal{M}) \right) \times \left( \bigtimes^r \Omega^1(\mathcal{M}) \right) \longrightarrow \mathcal{C}^\infty (\mathcal{M})$ \\
        oppure $\left(\bigtimes^k \mathcal{X}(\mathcal{M})\right) \times \left(\bigtimes^{r-1} \Omega^1(\mathcal{M})\right) \longrightarrow \mathcal{X}(\mathcal{M}),$ è indotta da un campo tensoriale misto di rango (k,r) \textbf{se e solo se} essa è \textbf{multilineare} su $\mathcal{C}^\infty (\mathcal{M})$.
    \end{lemma}
    Prestare attenzione a come si ripropone la questione degli indici tra un numero e il successivo ($r-1$ e $r$) a seconda della definizione della funzione e della natura del suo output: la prima mappa in una funzione reale liscia sulla varietà, la seconda mappa nello spazio vettoriale dato dalla sezione globale della varietà.
    %ORSO: lo spazio vettoriale X(M) è una sorta di covettore che mangia le funzioni lisce?
    
    Dati il tensore covariante di rango $k, \; \rho \in \mathcal{T}^k(\mathcal{N})$, e la mappa  liscia tra varietà $F : \mathcal{M} \longrightarrow \mathcal{N}$, si può effettuare il \emph{pull back} del tensore:
    \[
    (F^* \rho)_p (X_1, \dots, X_k) = \rho_{F(p)} (F_*X_1, \dots, F_*X_k)
    \]
    con $X_1, \dots, X_k$ campi vettoriali su $\mathcal{M}$.
    \begin{property} Si elencano alcune proprietà del pull back di funzioni applicate a tensori:
        \begin{tasks}[label=\arabic*.](2)
            \task $F^*$ è lineare in $\mathbb{R}$
            \task $F^*(f \rho) = (f \circ F) F^* \rho$
            \task $F^* (\rho \otimes \omega) = F^* \rho \otimes F^* \omega$
            \task $(G \circ F)^* = F^* \circ G^*$
        \end{tasks}
    \end{property}

    \section{Forme differenziali}
    \begin{definition}
        Dato $\Lambda^k \mathcal{M} = \bigsqcup_{p \in \mathcal{M}} \Lambda^k(T_p \mathcal{M})$, le sezioni globali $\Gamma(\Lambda^k \mathcal{M})$ sono dette \textbf{$k$-forme differenziali} o forme differenziali di rango k. L'insieme di queste ultime è detto $\Omega^k(\mathcal{M})$.
        Data una carta, si può scrivere
        \[
        \omega = \frac{1}{k!} \omega_{i_1, \dots, i_k} \odif{x^{i_1}} \wedge \cdots \wedge \odif{x^{i_k}} = \frac{1}{k!} \omega_{i_1, \dots, i_k} \left(\bigwedge_{j=1}^k \odif{x^{i_j}}\right)
        \]
    \end{definition}
    \begin{lemma}
        Sia $F : \mathcal{M} \longrightarrow \mathcal{N}$ una mappa liscia tra varietà e siano $(\mathbf{x})$ e $(\mathbf{y})$ le funzioni delle coordinate per $\mathcal{M}$ e $\mathcal{N}$. Allora
        \begin{equation}
        F^*\left(\omega_{i_1, \dots, i_k} \bigwedge_{j=1}^k \odif{y^{i_j}}\right) = \left( \omega_{i_1, \dots, i_k} \circ F\right) \bigwedge_{j=1}^k \odif{\left( y^{i_j} \circ F \right)}
        \end{equation}
        \begin{proof}
            \begin{align*}
                F^*\left(\omega_{i_1, \dots, i_k} \bigwedge_{j=1}^k \odif{y^{i_j}}\right) &= \left( \omega_{i_1, \dots, i_k} \circ F \right) F^* \left( \bigwedge_{j=1}^k \odif{y^{i_j}} \right) \qquad \footnotesize\text{(dalla proprietà 2.4.3.2)} \\
                \footnotesize\text{(dalla proprietà 2.4.3.3)} \qquad &= \left( \omega_{i_1, \dots, i_k} \circ F \right) \bigwedge_{j=1}^k F^*\odif{y^{i_j}} \\
                \footnotesize\text{(dalla proprietà 1.3.12.1)} \qquad &= \left( \omega_{i_1, \dots, i_k} \circ F \right) \bigwedge_{j=1}^k \odif{y^{i_j} \circ F}
            \end{align*}
        \end{proof}
    \end{lemma}

    Per una qualsiasi varità liscia, esiste un \emph{operatore differenziale}
    \[
    d : \Omega^k (\mathcal{M}) \longrightarrow \Omega^{k+1}(\mathcal{M})
    \]
    \begin{equation}
        \odif{\omega} \coloneqq \frac{1}{k!}\odif{\omega_{i_1, \dots, i_k}} \bigwedge_{j=1}^k \odif{x^{i_j}}
    \end{equation}
    o alternativamente
    \begin{equation}
        \odif{\omega} \coloneqq \frac{1}{(k+1)!} \left( \partial_{[j} \omega_{i_1, \dots, i_k]} \right) \odif{x^j} \wedge \left(\bigwedge_{l=1}^k \odif{x^{i_l}}\right)
    \end{equation}
    dove con $[j, i_1, \dots, i_k]$ si intende che si prendono gli indici completamente antisimmetrici.
    \begin{property} Si elencano le proprietà di questo operatore, detto anche \emph{derivata esterna}
        \begin{itemize}
            \item d è ben definita indipendentemente dalla scelta delle coordinate
            \item d è lineare su $\mathbb{R}$
            \item d$^2 = 0$
            \item dati $\omega \in \Omega^k(\mathcal{M})$ e $\rho \in \Omega^n(\mathcal{M})$, vale che $\odif{(\omega \wedge \rho)} = \odif{\omega} \wedge \rho + (-1)^k \omega \wedge \odif{\rho}$
        \end{itemize}
    \end{property}